% Full simulation-inclusive revision, August 2026.
% Rev. 14: valid-reference event-mean scalar; disclosed detector cache grid;
% Geant4 reference-list labels; synchronized code and output archive.
\documentclass[aps,pra,reprint,longbibliography,floatfix]{revtex4-2}

\usepackage[utf8]{inputenc}
\usepackage[T1]{fontenc}
\usepackage{amsmath,amssymb,bm}
\usepackage{booktabs}
\usepackage{graphicx}
\graphicspath{{out/}}
\usepackage{hyperref}
\usepackage{microtype}

% Shorthand notation used throughout the manuscript.
\newcommand{\epsM}{\ensuremath{\varepsilon_M}}
\newcommand{\thetazero}{\ensuremath{\theta_0}}
\newcommand{\thetaspace}{\ensuremath{\theta_{\mathrm{space}}}}
\newcommand{\thetarms}{\ensuremath{\theta_{\mathrm{rms}}}}
\newcommand{\thetacut}{\ensuremath{\theta_{\mathrm{cut}}}}
\newcommand{\dtheta}{\ensuremath{\Delta\theta}}
\newcommand{\XZero}{\ensuremath{X_0}}
\newcommand{\chic}{\ensuremath{\chi_c}}
\newcommand{\chia}{\ensuremath{\chi_a}}
\newcommand{\Inom}{\ensuremath{I_{\mathrm{nom}}}}
\newcommand{\Ip}{\ensuremath{I_p}}
\newcommand{\IQ}{\ensuremath{I_Q}}
\newcommand{\pathP}{\ensuremath{\mathcal{P}}}

\begin{document}
\title{The Highland Constant and the Divergent Second Moment:\\
Acceptance Dependence of Momentum-Informed Muon Scattering Tomography}
\author{Melora Wu}

\begin{abstract}
Momentum-informed muon scattering tomography often normalizes a quadratic scattering-angle weight by the Highland multiple-Coulomb-scattering width.  The Highland parameterization describes a Gaussian core, whereas a quadratic space-angle observable samples an acceptance-conditioned second moment of $\dtheta=\sqrt{\theta_x^2+\theta_y^2}$.  In the idealized small-angle Rutherford/Moli\`ere tail the two-dimensional angular density scales as $\Theta^{-4}$ and the corresponding magnitude density as $\Theta^{-3}$, so the second moment grows logarithmically when the angular response is extended to progressively larger angles.  Within the common-momentum, non-factorized radial Moli\`ere model we derive the exact reduced form $(1+\epsM)^2=RB\,\mu_2(\eta_{\rm cut};B)$ with $\eta_{\rm cut}=\thetacut/(\chic\sqrt{B})=k/\sqrt{2RB}$.  For a fixed path, $R$ and $B$ vary by less than $0.1\%$ over $1$--$6~\mathrm{GeV}/c$, so momentum enters almost entirely through reduced acceptance, while residual composition dependence remains through $RB$.  For the central $10~\mathrm{cm}$ Al $+15~\mathrm{cm}$ Cu path, the segmented $p(X)$ construction predicts incident-to-exit momentum changes from $-24.8\%$ at $1~\mathrm{GeV}/c$ to $-4.9\%$ at $6~\mathrm{GeV}/c$.  The self-consistent acceptance mismatch then rises from $1.66\%$ to $15.81\%$, whereas an upstream-momentum Highland denominator gives a much flatter $17.22$--$18.77\%$ mismatch, corresponding to an expected nominal quadratic weight of about $1.37$--$1.41$ under the reference model.  The detector study is implemented as a controlled model-internal PoCA simulation with analytically ordered ray segments evaluated on a disclosed cache grid, a $9\times9$ raster, explicit global-scale controls, truth/reconstructed path-class diagnostics, and a split-half map-noise estimate.  Across five $2\times10^6$-event realizations, quantitative map statistics restricted to occupancy-qualified voxels fully contained in the outer target give $\operatorname{RMS}(I_{\rm nom}-I_Q)=0.4616\pm0.0019$.  The optimal single global scale leaves $0.01870\pm0.00050$, while the central-path $p(X)$ correction leaves $0.01288\pm0.00021$, a further $31.1\pm0.8\%$ reduction of the post-scalar residual.  The segmented $p(X)$ construction remains a model extension rather than a transport theorem; independent single-slab Geant4 benchmarking is kept as a separate protocol.  Finite acceptance is therefore part of the normalization problem itself, while acceptance optimization and quantitative calibration remain distinct from task-specific image optimization.
\end{abstract}

\maketitle

\section{Introduction}
Muon scattering tomography exploits the dependence of multiple Coulomb scattering on atomic number and areal density to image large, dense objects non-invasively~\cite{miyadera2022}: high-$Z$ material at borders~\cite{schultz2004}, spent fuel in dry casks~\cite{bae2024,poulson2019}, and reactor structures~\cite{procureur2018}.  Reconstruction ranges from Point of Closest Approach (PoCA)~\cite{schultz2004} through maximum-likelihood EM~\cite{schultz2007} to more general trajectory and image-reconstruction methods that account for path uncertainty, energy loss, or curvature~\cite{chatzidakis2018,liu2019,ughade2025}.

Momentum information has already been incorporated into PoCA and related muon-scattering tomography in several forms, including momentum-dependent PoCA, categorized momentum--angle variables, and direct momentum rescaling of measured scattering angles~\cite{bae2022,bae2024ane,bae2024net,bae2024,yu2025}.  The estimator studied here is narrower: it normalizes each measured space angle by a Highland-predicted reference width and then squares the ratio [Eq.~\eqref{eq:weight}].  In the literature located for this audit, we did not identify a published treatment that derives the calibration of this specific Highland-normalized quadratic space-angle weight in terms of the acceptance-conditioned Moli\`ere second moment.  The novelty claimed here is therefore not momentum weighting itself, nor angular filtering, but the reduced-variable calibration identity for this quadratic estimator and its consequences for a reference-hypothesis normalization.

The distinction is already implicit in the Gaussian-core construction of Lynch and Dahl and in the modern PDG summary.  Their fitted coefficient depends on the central fraction of the projected distribution included in the Gaussian fit; the familiar $13.6~\mathrm{MeV}$ value corresponds to the central $98\%$~\cite{lynch1991,pdg2026}.  A quadratic space-angle weight, by contrast, samples a second moment.  Multiple-scattering theory separates a near-Gaussian soft component from a hard Rutherford component, and the naive second moment of the Rutherford asymptotics is logarithmically sensitive to the angular range~\cite{moliere1948,bethe1953,scott1963,bondarenco2016}.  Tomographic studies have likewise shown that large-angle events, angular cuts, and angle-dependent weights can materially change RMS- or PoCA-based image statistics~\cite{riggi2021,wen2023,wang2025}.  In the idealized small-angle Rutherford/Moli\`ere model used here, the two-dimensional tail falls as $\Theta^{-4}$ and, after the polar Jacobian, the probability density of the measured magnitude $\Theta=\sqrt{\theta_x^2+\theta_y^2}$ falls as $\Theta^{-3}$.  The corresponding second moment therefore grows logarithmically with the upper angular cut.  The formal divergence is a property of extending the small-angle model to arbitrarily large $\Theta$, not a claim that physical scattering is literally unbounded; Appendix~\ref{app:g4} defines the independent transport benchmark used to test the finite-acceptance prediction.

We denote the mismatch between the in-acceptance RMS and the Highland space-angle width by $\epsM$.  The central theoretical result of this work is that, within the common-momentum non-factorized radial Moli\`ere model, the mismatch can be written entirely in reduced variables as $(1+\epsM)^2=RB\mu_2(\eta_{\rm cut};B)$.  This identity separates the fixed-path momentum dependence---the reduced acceptance---from residual material dependence through $R$ and $B$.  We then test the limits of that separation: matched-$k$ and matched-$\eta_{\rm cut}$ comparisons across Al, Cu, Al+Cu, and Pb; explicit momentum loss along a thick reference path; event-level Al-only versus Cu-bearing path classes; a controlled spatial momentum-mixture gradient; and adaptive-cut retention for Pb- and Cu-dominated trajectories.

The numerical work remains deliberately controlled.  The fixed-$200~\mathrm{mrad}$ detector study uses the segmented $p(X)$ scattering construction both to generate the integrated target deflection and to construct the acceptance-matched reference denominator; it is therefore a model-internal reconstruction study rather than independent transport validation.  Appendix~\ref{app:g4} defines a separate single-material Geant4 benchmark protocol.  The stopping calculation uses the implemented Bethe collision model and neglects radiative loss.  The adaptive-cut diagnostic remains based on the constant-momentum radial optimum.  These qualifications matter for absolute quantitative calibration, but they do not affect the analytic distinction between a core width and the finite-acceptance second moment sampled by the estimator.

\section{Theoretical framework}\label{sec:framework}

\subsection{Highland core width}
For a muon traversing thickness $x$ of radiation length $\XZero$, the RMS projected deflection is the Highland formula as reparameterized by Lynch and Dahl~\cite{highland1975,lynch1991}, in the PDG form~\cite{pdg2026},
\begin{equation}\label{eq:highland}
\thetazero = \frac{13.6~\mathrm{MeV}}{\beta c p}\sqrt{\frac{x}{\XZero}}
\left[1+0.038\ln\!\left(\frac{x}{\XZero\beta^2}\right)\right],
\end{equation}
with $\beta=p/\sqrt{p^2+m_\mu^2}$.  We use
\begin{equation}
\thetaspace=\sqrt{2}\,\thetazero,
\end{equation}
so $\thetazero$ is the projected RMS and $\thetaspace$ is the corresponding two-dimensional Gaussian-core RMS.  Lynch and Dahl obtain $S_2=13.6~\mathrm{MeV}$ at a central fraction $F=98\%$, with $S_2=12.1-0.4\ln(1-F)$ over $0.7<F<0.997$~\cite{lynch1991}.  The coefficient therefore characterizes a specified central fraction rather than the full second moment.

\subsection{Radial truncated moments}\label{sec:moments}
The imaging observable is the non-negative space-angle magnitude
\begin{equation}\label{eq:Theta_def}
\Theta\equiv\dtheta=\sqrt{\theta_x^2+\theta_y^2}\ge0.
\end{equation}
Let $P(\Theta)$ be a rotationally symmetric two-dimensional angular density per unit angular area,
\begin{equation}\label{eq:Pnorm}
2\pi\int_0^\infty \Theta P(\Theta)\,d\Theta=1,
\end{equation}
and define the corresponding magnitude density
\begin{equation}\label{eq:hdef}
h(\Theta)=2\pi\Theta P(\Theta),\qquad \int_0^\infty h(\Theta)\,d\Theta=1.
\end{equation}
For an angular acceptance $\thetacut$,
\begin{equation}\label{eq:Fc}
F_c=\int_0^{\thetacut}h(\Theta)\,d\Theta,
\qquad
M_n(\thetacut)=\frac{\int_0^{\thetacut}\Theta^n h(\Theta)\,d\Theta}{F_c}.
\end{equation}
Thus
\begin{align}
\thetarms^2(\thetacut)&=M_2(\thetacut),\label{eq:rms_var_def}\\
\operatorname{Var}(\dtheta^2\mid\dtheta<\thetacut)&=M_4-M_2^2.
\end{align}
The signed projected density is useful as a cross-check, but it is not the probability measure used to define the tomography moments.

The hard radial cut of Eq.~\eqref{eq:Fc} is an idealization.  Experimental PoCA analyses already show that detector response, angle cuts, and angle-dependent event weighting can alter reconstructed signal distributions~\cite{riggi2021}.  A real detector applies a selection kernel $A(\bm{\theta},p,\pathP)$ built from tracker geometry, hit and track-quality requirements, reconstruction efficiency, and momentum resolution, and the sampled moment is then
\begin{equation}\label{eq:M2_kernel}
M_2[A]=\frac{\int \Theta^2 P(\bm{\theta})A(\bm{\theta},p,\pathP)\,d^2\theta}
{\int P(\bm{\theta})A(\bm{\theta},p,\pathP)\,d^2\theta},
\end{equation}
with Eq.~\eqref{eq:Fc} recovered for $A=\mathbb{1}[\Theta<\thetacut]$.  The conceptual point of this paper---that a quadratic weight samples an acceptance-dependent second moment rather than a core width---does not require radial symmetry or a sharp cut, but it does require the response to suppress sufficiently large angles that the accepted second moment exists.  For a Rutherford-like tail this requires the angularly averaged response $\overline A(\Theta)$ to satisfy $\int^\infty \overline A(\Theta)\,d\Theta/\Theta<\infty$; boundedness of $A$ alone is not sufficient.  The closed-form reduced identity of Sec.~\ref{sec:reduced} does require the hard-cut form, and all numerical work below uses it.

\subsection{Logarithmic second moment and quadratic fourth moment}\label{sec:tails}
A screened single-scatter Rutherford cross section, expressed per unit angular area in the two-dimensional plane, may be written~\cite{moliere1948,bethe1953,scott1963,bondarenco2016}
\begin{equation}\label{eq:rutherford2d}
P_{\rm R}(\Theta)=\frac{1}{\pi}\frac{\chic^2}{(\Theta^2+\chia^2)^2}
\longrightarrow \frac{\chic^2}{\pi\Theta^4}
\quad(\Theta\gg\chia).
\end{equation}
Equation~\eqref{eq:rutherford2d} is a single-scatter differential and is not normalized under Eq.~\eqref{eq:Pnorm}; its role is to fix the tail coefficient of the full multiple-scattering density $P_{\rm M}$ of Appendix~\ref{app:moliere}, which is normalized.  The corresponding magnitude tail is
\begin{equation}\label{eq:tail}
h(\Theta)\longrightarrow \frac{2\chic^2}{\Theta^3}.
\end{equation}
As $F_c\rightarrow1$,
\begin{equation}\label{eq:secondmoment}
M_2(\thetacut)\sim \mathrm{const}+2\chic^2\ln\thetacut,
\end{equation}
up to the dimensionful intercept implied by matching the tail to the core.  The corresponding fourth moment grows quadratically,
\begin{equation}\label{eq:m4}
M_4(\thetacut)\sim \mathrm{const}+\chic^2\thetacut^2.
\end{equation}
We define
\begin{equation}\label{eq:epsM}
\epsM(\thetacut)=\frac{\thetarms(\thetacut)-\thetaspace}{\thetaspace}.
\end{equation}
The logarithm is an asymptotic statement; finite-cut values must be obtained from the full accepted distribution.

\subsection{Non-factorized radial Moli\`ere evaluation}\label{sec:quadrature}
For a specified common momentum and material path, the primary numerical object is the rotationally symmetric Moli\`ere density $P_{\rm M}(\Theta)$ of Appendix~\ref{app:moliere}.  The accepted moments are evaluated directly from the radial probability measure,
\begin{align}
M_2(\thetacut)&=\frac{\int_0^{\thetacut}\Theta^3P_{\rm M}(\Theta)\,d\Theta}{\int_0^{\thetacut}\Theta P_{\rm M}(\Theta)\,d\Theta},\label{eq:quad}\\
M_4(\thetacut)&=\frac{\int_0^{\thetacut}\Theta^5P_{\rm M}(\Theta)\,d\Theta}{\int_0^{\thetacut}\Theta P_{\rm M}(\Theta)\,d\Theta}.\label{eq:M4_radial}
\end{align}
No Cartesian factorization $F(\theta_x)F(\theta_y)$ is used for the hard-scatter tail.  The radial expansion is truncated at $n\le2$.  Three-term Moli\`ere/Bethe expansions are accurate at the percent level in the core once $B$ is sufficiently large, but the $1/B$ expansion is asymptotic rather than uniformly convergent in the far tail~\cite{andreo1993,bondarenco2016}.  Appendix~\ref{app:checks} therefore quantifies the finite-cut numerical sensitivity directly rather than treating the truncation as uniformly controlled.

\subsection{Exact reduced-variable form}\label{sec:reduced}
Define the Moli\`ere scale and dimensionless series
\begin{equation}
s\equiv\chic\sqrt{B},\qquad \eta\equiv\Theta/s,
\qquad
Q_B(\eta)\equiv\sum_{n=0}^{2}B^{-n}\Phi^{(n)}(\eta),
\end{equation}
so that $P_{\rm M}(\Theta)=Q_B(\eta)/(2\pi s^2)$.  Define the dimensionless conditional second moment
\begin{equation}\label{eq:mu2}
\mu_2(\eta_c;B)=
\frac{\int_0^{\eta_c}\eta^3Q_B(\eta)\,d\eta}
{\int_0^{\eta_c}\eta Q_B(\eta)\,d\eta}.
\end{equation}
Then $M_2=s^2\mu_2$.  With
\begin{equation}\label{eq:Rdef}
R\equiv\frac{\chic^2}{\thetaspace^2},\qquad
k\equiv\frac{\thetacut}{\thetazero},
\end{equation}
the mismatch obeys the exact identity, within the common-$p$, $n\le2$ radial model,
\begin{equation}\label{eq:reduced_master}
\boxed{(1+\epsM)^2=RB\,\mu_2(\eta_{\rm cut};B)},
\end{equation}
with
\begin{equation}\label{eq:eta_cut}
\eta_{\rm cut}=\frac{\thetacut}{\chic\sqrt{B}}
=\frac{k}{\sqrt{2RB}}.
\end{equation}
Composition does not disappear from Eq.~\eqref{eq:reduced_master}: it remains through both $R$ and $B$.

The far-tail coefficient gives the asymptotic form
\begin{equation}\label{eq:eta1_asymptote}
\epsM\simeq\sqrt{1+2R\ln(\eta_{\rm cut}/\eta_1)}-1,
\end{equation}
where $\eta_1$ is a core-to-tail intercept, not a universal constant fixed by the Rutherford coefficient.  Appendix~\ref{app:checks} shows explicitly that the fitted slope approaches $2R$ while the fitted $\eta_1$ remains path dependent.

\subsection{Momentum invariance of $R$ and $B$}\label{sec:RBinvariance}
The momentum behavior of Eq.~\eqref{eq:reduced_master} can be made quantitative without simulation.  Combining Eqs.~\eqref{eq:highland} and~\eqref{eq:multi}, the $(p\beta)^{-2}$ factor cancels identically in $R$, leaving
\begin{equation}\label{eq:R_explicit}
R=\frac{0.157\sum_i Z_i(Z_i+1)X_i/A_i}
{2\,(13.6~\mathrm{MeV})^2\,(x/\XZero)_{\rm tot}\,L_{\rm H}^2},
\end{equation}
with $L_{\rm H}$ the bracketed Highland logarithm of Eq.~\eqref{eq:highland}.  The only surviving momentum dependence is the $\beta^2$ inside $L_{\rm H}$.  For the central Al+Cu path, $(x/\XZero)_{\rm tot}=11.57$ and $\sum_i Z_i(Z_i+1)X_i/A_i=2.02\times10^{3}~\mathrm{g\,cm^{-2}\,mol^{-1}}$, giving $R=0.0621$; between $1$ and $6~\mathrm{GeV}/c$, $\Delta\ln\beta^2=0.011$ and $R$ therefore varies by $0.08\%$.  Likewise $p^2$ cancels in the ratio $\chic^2/\chia^2$ that fixes Eq.~\eqref{eq:Bsolve}, whose right-hand side varies only through $\beta^{-2}$ and the $(Z\alpha/\beta)^2$ screening term; the two partially cancel, and $B=16.9$ varies by less than $0.1\%$ over the same range.  The product $RB$ is thus constant to better than $0.15\%$ across the four beam settings.

Two consequences follow.  First, at fixed physical cut the entire momentum dependence of $\epsM$ for a fixed path enters through $\eta_{\rm cut}$.  Since $\chic\propto(p\beta)^{-1}$, Eq.~\eqref{eq:eta_cut} gives $\eta_{\rm cut}\propto p\beta/\sqrt{B}$ exactly; the residual $B(p)$ dependence is below $0.1\%$ over the range studied, so $\eta_{\rm cut}\propto p\beta$ is an excellent approximation but not an identity.  Second, the reduced-variable collapse examined in Sec.~\ref{sec:collapse_results} is not an empirical finding but a consistency check on the quadrature: Eq.~\eqref{eq:reduced_master} forces the four points onto one curve at the $10^{-3}$ level once $RB$ is fixed.

\subsection{Momentum evolution through thick paths}\label{sec:eloss}
When momentum loss is appreciable, a single reduced momentum variable is insufficient.  Energy-loss-aware trajectory estimation has been studied in muon tomography~\cite{chatzidakis2018,ughade2025}; here the purpose is more limited, namely to construct a controlled mean $p(X)$ extension of the scattering normalization.  The mean energy is propagated according to
\begin{equation}\label{eq:dedx}
\frac{dE}{dX}=-S(E;\pathP),\qquad
p(X)c=\sqrt{E(X)^2-m_\mu^2c^4}.
\end{equation}
The scattering calculation then depends on the full profile $p(X)$.  In the numerical construction used below, the path is sliced until the fractional change of $p\beta$ within a slice is below $1\%$; $\chic^2$ is accumulated with the local $(p_j\beta_j)^{-2}$ factor, the screening scale is combined logarithmically, and $B$ is solved once for the accumulated path.  A corresponding $p(X)$ Highland-core construction integrates the local $(p\beta)^{-2}$ scattering strength and applies one global logarithmic correction.  These varying-$p$ combinations reduce to the standard common-$p$ formulas in the weak-loss limit, but they are modeling choices rather than a new closed-form Highland theorem; Appendix~\ref{app:pofx} states the construction and its limitations explicitly.

We distinguish two mismatches.  The self-consistent quantity
\begin{equation}\label{eq:eps_profile}
\epsM[p(X)]\equiv\frac{\thetarms[p(X)]}{\thetaspace[p(X)]}-1
\end{equation}
uses the same momentum profile in numerator and core reference.  The deployed upstream-tagged mismatch instead compares $\thetarms[p(X)]$ with a Highland denominator evaluated at the incident tagged momentum.  The latter is the relevant quantity when the detector measures momentum upstream while scattering occurs after substantial energy loss.

\subsection{Momentum-informed PoCA normalization}\label{sec:poca}
PoCA~\cite{schultz2004} places the vertex at the midpoint of the shortest segment between extrapolated upstream and downstream tracks.  Momentum-dependent PoCA and momentum--angle imaging variables provide the closest published precedents for the present event-level normalization~\cite{bae2022,bae2024ane,bae2024net,bae2024,yu2025}.  The nominal event weight is
\begin{equation}\label{eq:weight}
w_{\rm nom}=\left(\frac{\dtheta}{\thetaspace(p_{\rm in},\pathP_{\rm reco}\mid\mathrm{ref})}\right)^2,
\end{equation}
where the reference geometry contains Al and Cu but excludes the Pb inclusion.  Under the reference hypothesis, the expected nominal weight is not generally unity because the numerator samples $M_2(\thetacut)$ while the denominator is a core width.

The acceptance-matched target is
\begin{equation}\label{eq:wQ}
w_Q=\left(\frac{\dtheta}{\thetarms(\thetacut,p,\pathP\mid\mathrm{ref})}\right)^2,
\end{equation}
for which
\begin{equation}\label{eq:EwQ}
\mathbb{E}[w_Q\mid \dtheta<\thetacut,p,\pathP,\mathrm{ref}]=1
\end{equation}
at the idealized true-parameter level.  A detector-level implementation requires conditioning on measured momentum and reconstructed path, including their resolution.  The production image $\IQ$ below is therefore a reconstructed-variable plug-in implementation of Eq.~\eqref{eq:wQ}, not a claim of exact response-conditioned unit expectation.

The denominator of Eq.~\eqref{eq:wQ} also requires a material hypothesis along the event path, since $R$, $B$, and $p(X)$ are all path-composition dependent.  The present study is therefore a reference-hypothesis calibration: the Al and Cu geometry is taken as known and the Pb inclusion is the sought deviation, which is the natural setting for cask or container monitoring against a nominal loading.  It is not directly a normalization scheme for unconstrained tomography, where composition is itself the unknown.  Applying an acceptance-matched denominator in that setting would require an iterative scheme, a deliberately reference-independent approximation, or a denominator marginalized over plausible material paths; none of these is developed here.

\subsection{Acceptance efficiency is not a universal imaging objective}\label{sec:optcut}
At fixed incident flux, a useful single-material moment-efficiency measure is
\begin{equation}\label{eq:eta}
\eta(\thetacut)=\sqrt{F_c(\thetacut)}\,
\frac{M_2(\thetacut)}{\sqrt{M_4(\thetacut)-M_2(\thetacut)^2}}.
\end{equation}
Because $M_2$ grows only logarithmically while $M_4$ grows quadratically, $\eta$ has a finite optimum for a specified model and path.  That optimum changes the estimand by changing the accepted moment and need not optimize a tomographic task such as Pb--Cu discrimination.  Prior PoCA studies likewise find material changes when angle cuts or large-angle weighting are altered~\cite{riggi2021,wen2023,wang2025}.  Acceptance optimization must therefore be performed together with an acceptance-matched denominator and against the actual imaging objective.

\section{Controlled numerical study and results}\label{sec:results}

\subsection{Fixed-path momentum collapse}\label{sec:collapse_results}
Figure~\ref{fig:collapse} evaluates Eq.~\eqref{eq:reduced_master} for the central Al+Cu reference path.  The four physical momentum settings, evaluated at the fixed $200~\mathrm{mrad}$ cut, lie on the single reduced curve to within the $0.15\%$ invariance of $RB$ established in Sec.~\ref{sec:RBinvariance}.  This is a closure check on the radial quadrature rather than a discovery: momentum changes only how far into the radial tail the detector acceptance extends, and $\eta_{\rm cut}$ is the coordinate along which the physical cut moves.

\begin{figure}[t]
\centering
\includegraphics[width=\columnwidth]{figs/theory_collapse_eta}
\caption{Common-momentum reduced-variable collapse for the central Al+Cu reference path.  The line is Eq.~\eqref{eq:reduced_master}; points are the four beam momenta at a fixed $200~\mathrm{mrad}$ physical cut.  The collapse is forced by the momentum invariance of $RB$ (Sec.~\ref{sec:RBinvariance}) and is shown as a numerical closure test.}
\label{fig:collapse}
\end{figure}

The identity is not a statement of composition universality.  Figure~\ref{fig:composition} compares Al $25~\mathrm{cm}$, Cu $15~\mathrm{cm}$, Al+Cu, and Pb $15~\mathrm{cm}$ at matched $k$ and matched $\eta_{\rm cut}$.  Material differences remain in both comparisons: at matched $\eta_{\rm cut}=10$, $\epsM$ is $13.2\%$ for Al+Cu and $19.2\%$ for Pb.  Equation~\eqref{eq:reduced_master} localizes that difference to the prefactor.  Using $R=0.0621$, $B=16.9$ for Al+Cu and $R=0.0702$, $B=16.5$ for Pb gives $\mu_2(10;B)=1.224$ and $1.227$ respectively, so the reduced conditional moment is common to the two paths at the few-per-mille level and essentially all of the residual composition dependence is carried by $RB$.  The practical statement is therefore sharper than "matching $\eta_{\rm cut}$ is insufficient": $\epsM$ at matched reduced acceptance is predicted by $\sqrt{RB\,\mu_2}$ with a nearly path-independent $\mu_2$, so a per-path correction requires only the two Moli\`ere parameters and no re-evaluation of the radial integral.

\begin{figure*}[t]
\centering
\includegraphics[width=0.48\textwidth]{figs/theory_composition_matched_k}\hfill
\includegraphics[width=0.48\textwidth]{figs/theory_composition_matched_eta}
\caption{Composition dependence after controlling the acceptance scale.  (a) $\epsM$ at matched $k=\thetacut/\thetazero$.  (b) The same paths at matched reduced acceptance $\eta_{\rm cut}$.  Matching $\eta_{\rm cut}$ removes the momentum/scale effect but not the material dependence, which enters through the prefactor $RB$.}
\label{fig:composition}
\end{figure*}

\subsection{Energy loss breaks the single-scale picture}\label{sec:eloss_results}
The central axial reference path contains $10~\mathrm{cm}$ of Al and $15~\mathrm{cm}$ of Cu, a thick reference path.  Figure~\ref{fig:eloss} compares the common-$p$ picture with the segmented $p(X)$ construction.  The predicted momentum change ranges from about $-25\%$ at $1~\mathrm{GeV}/c$ to about $-5\%$ at $6~\mathrm{GeV}/c$.

The two mismatches of Sec.~\ref{sec:eloss} behave very differently.  With the production-default $d\chi_c^2$ screening-log weighting, the self-consistent mismatch is $1.66$, $9.11$, $12.71$, and $15.81\%$ at $1$, $2$, $3.5$, and $6~\mathrm{GeV}/c$.  When the denominator instead uses the upstream tagged momentum, the effective mismatch is $17.38$, $17.22$, $17.54$, and $18.77\%$: a large offset with only a $1.56$ percentage-point spread.  The corresponding reference-hypothesis expectation of the nominal quadratic weight, $(1+\epsM)^2$, is $1.378$, $1.374$, $1.382$, and $1.411$.  Energy loss therefore changes both the magnitude and the momentum ordering of the deployed normalization error, and in the deployed case it largely removes the momentum dependence that the self-consistent quantity displays.

\begin{figure*}[t]
\centering
\includegraphics[width=\textwidth]{figs/theory_energy_loss_effect}
\caption{Energy-loss effect for the central Al+Cu path in the segmented $p(X)$ construction.  (a) Fractional incident-to-exit momentum change.  (b) Self-consistent finite-acceptance mismatch (both numerator and core reference use $p(X)$) and the upstream-tagged mismatch relevant when the Highland denominator uses the incident measured momentum.  The upstream-tagged curve is nearly flat across the beam.}
\label{fig:eloss}
\end{figure*}

These values are conditional on the present stopping and varying-$p$ construction.  The stopping model includes Bethe collision loss with a density-effect correction but not radiative loss, and the segmented degrading-path construction has not yet been validated against an independent layered transport run.  Figure~\ref{fig:eloss} is therefore used to establish the size and direction of the modeled effect, not as a final transport uncertainty budget.

\subsection{Production PoCA configuration and scope}\label{sec:simsetup}
The imaging target is a $25~\mathrm{cm}$ Al outer cube containing a $15~\mathrm{cm}$ Cu cube and a Pb cylinder of radius $2~\mathrm{cm}$ and height $15~\mathrm{cm}$ centered at $(x,y)=(3,2)~\mathrm{cm}$.  The denominator reference contains Al and Cu only: in the reference geometry the Pb volume is replaced by Cu.  Six tracking planes are located at $z=(-120,-90,-45,-15,25,65)~\mathrm{cm}$ with $200~\mu\mathrm{m}$ single-hit resolution.  The upstream momentum tag is represented by a point dipole with $B=1~\mathrm{T}$ and $L_{\rm eff}=0.30~\mathrm{m}$.  The nominal settings are $1$, $2$, $3.5$, and $6~\mathrm{GeV}/c$, each with a $1\%$ Gaussian distribution of true incident momentum and $2~\mathrm{mrad}$ beam divergence.

A $1~\mathrm{cm}$ Gaussian spot is rastered on a $9\times9$ grid with node centers spanning $\pm11~\mathrm{cm}$ in both transverse coordinates.  The beamline is re-steered independently at each nominal momentum so that the nominal dipole displacement does not generate a setting-dependent target offset.  Each production realization is defined to contain $5\times10^5$ generated muons per momentum setting, or $2\times10^6$ events total, with equal exposure across raster nodes within each setting.

The production generator is deliberately model internal.  For each event the incident ray is intersected analytically with the nested target and stored as the ordered five-segment path $[\mathrm{Al}_{\rm up},\mathrm{Cu}_{\rm up},\mathrm{Pb},\mathrm{Cu}_{\rm down},\mathrm{Al}_{\rm down}]$.  Scattering is sampled from the $n\le2$ non-factorized radial Moli\`ere distribution after the segmented $p(X)$ accumulation of Appendix~\ref{app:pofx}.  To cache this calculation, incident momentum is rounded to $0.010~\mathrm{GeV}/c$, each ordered segment's areal density to $0.25~\mathrm{g\,cm^{-2}}$, and the angular cut to $0.002~\mathrm{rad}$; segment order itself is retained.  The resulting integrated angular deflection is represented as a single equivalent kink at the midpoint of the traversed outer-target interval, after which downstream tracker hits are propagated and smeared.  There is no distributed target scattering in the production generator, and tracker/spectrometer material is not modeled.  The imaging study is therefore a controlled reconstruction and normalization experiment within the stated scattering model, not an independent transport validation.

The measured momentum is reconstructed from the difference between the two upstream track-arm slopes, $p_{\rm meas}=0.3BL/|\Delta\theta_{\rm bend}|$.  The hit geometry alone gives a bend-angle resolution of about $1.33~\mathrm{mrad}$, corresponding to a leading geometric fractional resolution of roughly $1.5$, $3.0$, $5.2$, and $8.9\%$ over the four settings.  The $1\%$ beam bite is a distribution of true momentum, not part of this detector resolution.  Because the estimator inverts a noisy bend angle and the normalization is nonlinear in $p_{\rm meas}$, no assumption of an unbiased Gaussian fractional momentum error is made; momentum-response bias is instead measurable directly from the simulated event table.

The fixed-cut analysis uses $\thetacut=200~\mathrm{mrad}$.  Images are accumulated on a $50^3$ grid covering $[-15,15]~\mathrm{cm}$ in each coordinate, corresponding to $0.6~\mathrm{cm}$ voxels.  Quantitative map statistics require both at least 20 entries and full containment of the voxel inside the $25~\mathrm{cm}$ outer Al cube: for a voxel center $(x_v,y_v,z_v)$, $|x_v|+0.3$, $|y_v|+0.3$, and $|z_v|+0.3$ must each be no larger than $12.5~\mathrm{cm}$.  This fiducial condition removes boundary-straddling voxels in which a grazing reconstructed reference path can produce an anomalously small normalization denominator.  The Pb ROI is the radius-$2~\mathrm{cm}$, $|z|\le7.5~\mathrm{cm}$ cylinder centered on the Pb inclusion.  The Cu comparison ROI is a radius-$3.75~\mathrm{cm}$, $|z|\le7.5~\mathrm{cm}$ cylinder centered on the target axis with the Pb ROI removed.  The reported detector metrics are
\begin{equation}
{\rm SNR}_{\rm Pb}=\frac{\overline I_{\rm Pb}}{s_{\rm Pb}},\qquad
{\rm CNR}=\frac{\overline I_{\rm Pb}-\overline I_{\rm Cu}}{s_{\rm Cu}},
\end{equation}
where the bars and $s$ denote the voxel mean and sample standard deviation within the corresponding occupancy-masked ROI.

\subsection{Implemented normalization controls}\label{sec:artifactresults}
Five image constructions are retained in the analysis.  $\Inom$ uses the reconstructed reference-path Highland core width evaluated at $p_{\rm meas}$.  $\Ip$ multiplies that event-specific Highland core width by the central-path $p(X)$ mismatch factor evaluated at the same measured momentum; it therefore removes the dominant axial acceptance/energy-loss correction while preserving the event-specific Highland path length.  $\IQ$ uses the reconstructed ordered reference path in the full segmented-$p(X)$ truncated-RMS denominator.  $I_{\rm ideal}$ uses true incident momentum, true reference path, and the true generated angle and is retained only as an internal closure.  Finally, $I_{\rm const}$ applies the unweighted event-count mean mismatch over accepted events with a defined, nonzero reconstructed reference path as a pure scale-null control.

A second scalar control is optimized directly at the image level.  For the common valid voxel set, the code evaluates
\begin{equation}
 c_* = \arg\min_c \sum_v\left(\frac{I_{{\rm nom},v}}{c}-I_{Q,v}\right)^2
 =\frac{\sum_v I_{{\rm nom},v}^2}{\sum_v I_{{\rm nom},v}I_{Q,v}},
\end{equation}
and reports the RMS of $I_{\rm nom}/c_*-I_Q$ alongside the event-mean scalar control and the $p(X)$ correction.  This separates a nearly global calibration offset from improvements that require momentum- or path-dependent structure.

Path residuals are evaluated separately using both truth and reconstructed reference-path classifications.  In addition, the code performs an independent random split of the accepted events and estimates the full-sample map-noise contribution from the difference of the two half-sample residual maps.  For half counts $n_0$ and $n_1$ in a voxel, the half difference is scaled by $\sqrt{n_0n_1}/(n_0+n_1)$ to estimate the standard deviation of the corresponding full-sample mean.  A quadrature-subtracted residual is reported only as a diagnostic of whether an observed path-class RMS is resolved above the sampling floor.

The fourth moment is used directly for occupancy assessment.  For the central Al+Cu $p(X)$ calibration, $M_4/M_2^2$ rises from $1.99$ at $1~\mathrm{GeV}/c$ to $11.05$ at $6~\mathrm{GeV}/c$, corresponding to $\sigma(w_Q)$ from $0.99$ to $3.17$.  At the permissive $N=20$ voxel threshold, the latter gives a standard error of about $0.71$ on a unit expected weight.  The split-half diagnostic is therefore required before a short-path map-RMS excess is interpreted as physical path structure.

\subsection{Fiducial image-normalization results}\label{sec:fiducial_results}
The five production realizations are statistically independent but use the same exposure and reconstruction.  Over the common fiducial and occupancy-qualified voxel population, the mean nominal difference is
\begin{equation}
\operatorname{RMS}(I_{\rm nom}-I_Q)=0.46158\pm0.00188,
\end{equation}
where the uncertainty is the seed-to-seed standard deviation.  The valid-reference event-mean scalar control leaves $0.01872\pm0.00051$.  Consistently, the least-squares global scale is $c_*=1.38515\pm0.00037$ and leaves
\begin{equation}
\operatorname{RMS}(I_{\rm nom}/c_*-I_Q)=0.018699\pm0.000499,
\end{equation}
a $95.95\pm0.10\%$ reduction from the nominal difference.  The central-path $p(X)$ correction leaves
\begin{equation}
\operatorname{RMS}(I_p-I_Q)=0.012878\pm0.000205,
\end{equation}
a $97.210\pm0.041\%$ reduction from nominal and, more diagnostically, a further $31.1\pm0.8\%$ reduction relative to the residual remaining after the best global scalar.  Thus most of the original normalization error is a common multiplicative scale, but a smaller reproducible component is removed only by the momentum-dependent $p(X)$ correction.

\begin{figure}[t]
\centering
\includegraphics[width=\columnwidth]{figs/production_fiducial_summary}
\caption{Five-seed fiducial image-RMS comparison.  Error bars are seed-to-seed standard deviations.  The optimized global scale removes most of the nominal difference, while the central-path $p(X)$ correction reduces the post-scalar residual by a further $31.1\%$.  The vertical axis is logarithmic.}
\label{fig:production_fiducial}
\end{figure}

A representative central-$z$ difference map is shown in Fig.~\ref{fig:diffmaps}.  The fiducial mask removes the partially outside outer voxel ring that otherwise produced rare grazing-path denominator instabilities.  The remaining $I_p-I_Q$ structure is concentrated mainly near the transition between Cu-bearing and Al-only reference paths rather than in the central Pb/Cu ROIs.

\begin{figure*}[t]
\centering
\includegraphics[width=\textwidth]{figs/seed0_difference_maps}
\caption{Representative fiducial central-slice difference maps for one of the five production seeds.  Left: nominal Highland normalization minus the acceptance-matched image.  Right: residual after the central-path $p(X)$ correction.  Voxels with fewer than 20 entries or not fully contained in the outer target are masked.}
\label{fig:diffmaps}
\end{figure*}

The earlier apparent large Al-only/Cu-bearing RMS ratio is not retained as a physical result.  Before the fiducial boundary correction, the reconstructed Al-only residual was strongly seed dependent and its split-half noise estimate was comparable to the observed map RMS, whereas the Cu-bearing residual was much more stable.  The revised code applies the same fiducial mask to the path-class and split-half analyses, and no path-class ratio is claimed here without a resolved post-mask excess above that noise floor.  Truth/reconstructed Al-only classification itself is stable: the event-level migration fraction over the five runs is $0.259\pm0.004\%$.

\subsection{Five-seed ROI metrics}\label{sec:paired}
The absolute image metrics averaged over the five realizations are $\mathrm{SNR}_{\rm Pb}=2.5389\pm0.0754$ and $\mathrm{CNR}=2.9645\pm0.0806$ for $I_{\rm nom}$, compared with $2.5121\pm0.0741$ and $2.9309\pm0.0791$ for $I_Q$.  The paired differences are
\begin{align}
\Delta\mathrm{SNR}_{\rm nom-Q}&=0.02684\pm0.00170,\\
\Delta\mathrm{CNR}_{\rm nom-Q}&=0.03357\pm0.00174.
\end{align}
After the central-path $p(X)$ correction, the corresponding residual differences are only $0.000800\pm0.000065$ and $0.000972\pm0.000156$.  A global multiplication does not change either metric by definition, so the scalar controls are omitted from Fig.~\ref{fig:paired}.

\begin{figure*}[t]
\centering
\includegraphics[width=0.86\textwidth]{figs/paired_seed_summary}
\caption{Paired five-seed changes in Pb SNR and Pb--Cu CNR relative to the acceptance-matched image $I_Q$.  Error bars are seed-to-seed standard deviations.  The central-path $p(X)$ correction reproduces the $I_Q$ ROI metrics to the $10^{-3}$ level in this controlled model.}
\label{fig:paired}
\end{figure*}

\subsection{Adaptive-cut diagnostic}\label{sec:adaptive}
The adaptive-cut diagnostic uses the constant-momentum radial optimum of Eq.~\eqref{eq:eta}; it is not the $p(X)$-aware fixed-cut calibration.  Across five seeds, truth-tagged Pb-crossing retention is $42.56\pm0.27\%$, compared with $72.44\pm0.03\%$ for Cu-only events.  Reconstructed ROIs give the same ordering, $37.52\pm0.31\%$ in the Pb ROI versus $69.89\pm0.20\%$ in the Cu ROI.  Thus the variance-based single-material cut preferentially removes the high-scattering population that carries Pb--Cu contrast.

\begin{figure*}[t]
\centering
\includegraphics[width=0.86\textwidth]{figs/adaptive_retention_ensemble}
\caption{Adaptive-cut retention averaged over five production seeds.  Error bars are seed-to-seed standard deviations.  The selection retains substantially fewer Pb-crossing/Pb-ROI events than Cu-only/Cu-ROI events.}
\label{fig:adaptive_retention}
\end{figure*}

\subsection{Controlled spatial momentum-mixture gradient}\label{sec:gradient}
The gradient exposure uses a reference-only Al+Cu target and fixes the total incident fluence per raster cell while varying the mixture of the four momentum settings across $x$.  With the same fiducial mask used for the production images, the unweighted normalization-field predictor has correlation $r=0.101$ with the observed difference map, and a best one-amplitude fit leaves $59.9\%$ of the observed uncentered RMS.  The eventwise $w_Q$-weighted construction gives unit amplitude and correlation with residual at numerical roundoff, as required by the algebraic identity $w_{\rm nom}-w_Q=w_Q[(1+\epsilon)^2-1]$.  Because the upstream-tagged mismatch spans only $1.56$ percentage points across the four settings, the imposed momentum-composition gradient has little dynamic range; the low correlation is therefore treated as inconclusive rather than evidence against a local normalization mechanism.

\section{Discussion}\label{sec:discussion}
The central analytic result is the reduced identity in Eq.~\eqref{eq:reduced_master}, together with the invariance argument of Sec.~\ref{sec:RBinvariance}.  Because $R$ and $B$ vary by less than $0.1\%$ for a fixed path over the momentum range studied, the common-momentum collapse is predictive rather than fitted and most of the fixed-cut momentum dependence is carried by $\eta_{\rm cut}$.  The same result prevents an over-interpretation of that collapse: material composition remains through $RB$.  At matched $\eta_{\rm cut}=10$, the Al+Cu and Pb examples give $\epsM=13.2\%$ and $19.2\%$, while their reduced conditional moments differ only at the few-per-mille level.  A universal momentum-only correction is therefore not exact for all reference paths.

Energy loss changes the problem more substantially than a small perturbation of reduced acceptance.  In the present segmented construction, the central Al+Cu path loses $24.8\%$ of its momentum at $1~\mathrm{GeV}/c$ and $4.9\%$ at $6~\mathrm{GeV}/c$.  The self-consistent acceptance mismatch spans $1.66$--$15.81\%$, but the upstream-tagged mismatch stays in the much narrower $17.22$--$18.77\%$ range.  Consequently the nominal quadratic weight under the reference model is expected to be biased upward by a nearly common factor of about $1.37$--$1.41$.  The five-seed image result confirms that most of the effect is scalar: $c_*=1.38515\pm0.00037$ reduces the fiducial RMS by $95.95\pm0.10\%$.  The central-path $p(X)$ correction reduces the remaining post-scalar residual by another $31.1\pm0.8\%$, demonstrating a smaller but reproducible non-scalar component within the controlled model.

The varying-$p$ screening average is not unique.  The production default weights the local screening logarithm by $d\chi_c^2$, while the alternative retains the common-$p$ serial $Z(Z+1)X/A$ weights slice by slice.  The two rules are algebraically identical at constant momentum.  On the central Al+Cu path their difference in the deployed upstream-tagged mismatch is $0.062$, $0.016$, $0.005$, and $0.002$ percentage point over the four settings.  This small spread is useful as an internal construction diagnostic, but it is not a substitute for independent transport validation of the one-$B$ degrading-path representation.

The detector simulation should be interpreted at the level at which it is constructed.  It includes beam divergence, finite spot size, a momentum bite, hit smearing, explicit bend-angle reconstruction, exact straight-ray intersections with the nested target, and PoCA reconstruction.  It does not include distributed target scattering, tracker material, magnet material, or spectrometer multiple scattering.  The total target deflection is sampled from the same segmented radial model that supplies the $I_Q$ denominator and is represented as one midpoint kink.  Thus the production study measures how the proposed normalization behaves under reconstruction and path uncertainty within the model; it does not validate the scattering model itself.

The revised path diagnostics also separate two questions that were previously conflated.  Truth and reconstructed Al-only/Cu-bearing classifications agree at the $99.74\%$ level, but that does not by itself establish a physical path-class RMS difference.  The original boundary-inclusive Al-only residual was seed unstable and split-half-noise dominated, which motivated the fully contained-voxel fiducial mask.  We therefore withdraw the earlier factor-$4.6$ path-class interpretation and require any future path-class claim to remain resolved above the split-half floor after the same fiducial selection.  The strong fourth-moment growth makes that requirement especially important for short, high-momentum paths.

Acceptance optimization remains a separate question from calibration.  The single-material criterion of Eq.~\eqref{eq:eta} changes both the estimator variance and the accepted material populations.  A cut that is favorable for estimating a single-material moment can preferentially suppress the high-scattering events that carry high-$Z$ contrast.  Quantitative normalization should therefore be matched to the deployed acceptance, while the acceptance itself should be optimized against the actual imaging task.

Several limitations remain.  The segmented $p(X)$ Moli\`ere accumulation and the $p(X)$ Highland-core construction are controlled extrapolations of common-$p$ formulas, not independently validated transport laws.  Collision stopping power is included but radiative loss is omitted, and the transcribed Sternheimer density-effect inputs require verification against primary tables before a precision stopping-power uncertainty is quoted.  The radial Moli\`ere expansion is truncated at $n\le2$; published analyses of the Moli\`ere series support percent-level core accuracy in the relevant large-$B$ regime but also emphasize its asymptotic tail behavior~\cite{andreo1993,bondarenco2016}.  In the present finite-cut calculation, the $n\le1$ versus $n\le2$ shift in $\epsM$ remains below $0.07$ percentage point for the four fixed-cut settings.  Finally, $I_Q$ is a reconstructed-variable plug-in rather than a fully response-conditioned detector calibration.

The separate Geant4 code supplied with this work is therefore treated as a benchmark protocol rather than as an extension of the production generator.  It uses single-material slabs, explicit run seeds, and runtime process-list reporting; the Python comparator evaluates the manuscript sign convention $\thetarms^{\rm model}/\thetarms^{\rm G4}-1$, the associated quadratic-weight bias, a median-based core diagnostic, and finite reduced-angle bands of the second-moment numerator.  Numerical transport claims should be inserted only from fresh angle dumps generated with a documented Geant4 version and physics configuration.

\section{Conclusions}
A Highland-normalized quadratic space-angle estimator compares two mathematically different quantities: a Gaussian-core width and an acceptance-conditioned second moment.  In the idealized small-angle Rutherford/Moli\`ere model, the latter grows logarithmically as the angular response is extended because the magnitude tail scales as $\Theta^{-3}$.  Within the common-momentum non-factorized radial model,
\begin{equation*}
(1+\epsM)^2=RB\mu_2(\eta_{\rm cut};B),\qquad
\eta_{\rm cut}=\frac{k}{\sqrt{2RB}}.
\end{equation*}
For a fixed path, $R$ and $B$ are nearly momentum invariant over the range studied, so reduced acceptance carries almost all of the common-$p$ momentum dependence while composition remains through $RB$.

For thick paths, energy loss prevents a single incident-momentum scale from describing both the core and the accepted second moment.  In the central Al+Cu example, the segmented construction predicts a large but weakly momentum-dependent upstream-tagged normalization mismatch.  In the five-seed fiducial image analysis, the best global scalar removes $95.95\pm0.10\%$ of the nominal map-RMS difference, while the central-path $p(X)$ correction removes $97.210\pm0.041\%$ and lowers the post-scalar residual by a further $31.1\pm0.8\%$.  The detector result is therefore not evidence for a large momentum-dependent artifact; it resolves a dominant global scale bias plus a smaller reproducible correction beyond that scale.

The controlled PoCA simulation uses analytically ordered ray segments evaluated on the disclosed cache grid, a $9\times9$ raster that explicitly samples Al-only reference paths, reconstructed momentum from the upstream bend, absolute ROI metrics, and a fully contained-voxel fiducial mask.  The prior large Al-only/Cu-bearing residual ratio is not retained because its boundary-inclusive short-path component was not resolved above the split-half noise floor.  Because the production generator and $I_Q$ share the same segmented scattering model, the detector results test reconstruction and normalization within that model rather than transport accuracy.  Independent transport benchmarking remains necessary for an absolute material scale.

The practical conclusion is therefore limited but robust: quantitative use of a quadratic scattering weight requires a denominator matched to the accepted second moment rather than only to the Highland core; appreciable energy loss must be propagated consistently; and angular acceptance must be optimized against the imaging task rather than treated as a universal efficiency constant.

\appendix

\section{Moli\`ere scattering parameters}\label{app:moliere}
Following the radial Moli\`ere/Bethe expansion~\cite{moliere1948,bethe1953} and the parameter conventions summarized by Lynch and Dahl~\cite{lynch1991}, with $p$ in $\mathrm{MeV}/c$ and areal density $X$ in $\mathrm{g\,cm^{-2}}$, define
\begin{equation}\label{eq:scale}
s\equiv\chic\sqrt{B},\qquad \eta\equiv\Theta/s.
\end{equation}
For the non-factorized two-dimensional distribution, the $n\le2$ radial Moli\`ere expansion is
\begin{equation}\label{eq:radial_moliere}
P_{\rm M}(\Theta)=\frac{1}{2\pi s^2}
\sum_{n=0}^{2}\frac{1}{B^n}\Phi^{(n)}(\eta),
\end{equation}
with Hankel-generating functions
\begin{equation}\label{eq:hankel_gen}
\Phi^{(n)}(\eta)=\frac{1}{n!}\int_0^\infty uJ_0(\eta u)e^{-u^2/4}
\left[\frac{u^2}{4}\ln\!\left(\frac{u^2}{4}\right)\right]^n du.
\end{equation}
For $n=0$, $\Phi^{(0)}(\eta)=2e^{-\eta^2}$ and
\begin{equation}
P_{\rm M}^{(0)}(\Theta)=\frac{e^{-\Theta^2/s^2}}{\pi s^2}.
\end{equation}
The normalized magnitude density is
\begin{equation}\label{eq:h_moliere}
h_{\rm M}(\Theta)=2\pi\Theta P_{\rm M}(\Theta),
\end{equation}
and the Rutherford regime approaches
\begin{equation}\label{eq:radial_tailnorm}
P_{\rm M}(\Theta)\longrightarrow\frac{\chic^2}{\pi\Theta^4},\qquad
h_{\rm M}(\Theta)\longrightarrow\frac{2\chic^2}{\Theta^3}.
\end{equation}

For comparison, the signed projected marginal is
\begin{equation}
F(\theta_x)=\int_{-\infty}^{\infty}P_{\rm M}\!\left(\sqrt{\theta_x^2+\theta_y^2}\right)d\theta_y,
\end{equation}
and may be represented by
\begin{equation}
F(\theta_{\rm plane})=\frac{1}{s}\sum_{n=0}^{2}\frac{1}{B^n}f_p^{(n)}(\eta_p),
\qquad \eta_p=\theta_{\rm plane}/s,
\end{equation}
with
\begin{equation}
f_p^{(n)}(\eta_p)=\frac{1}{\pi n!}\int_0^\infty
\cos(\eta_pu)e^{-u^2/4}
\left[\frac{u^2}{4}\ln\!\left(\frac{u^2}{4}\right)\right]^n du.
\end{equation}
The projected representation is retained only as an implementation cross-check; the tomography moments are always formed from $h_{\rm M}(\Theta)$.

For a single material,
\begin{align}
\chic^2&=\frac{0.157\,Z(Z+1)X/A}{(p\beta)^2},\label{eq:chic}\\
\chia^2&=\frac{2.007\times10^{-5}Z^{2/3}[1+3.34(Z\alpha/\beta)^2]}{p^2},\label{eq:chia}
\end{align}
with $\alpha=1/137.036$ and charge $z=1$.  The Moli\`ere parameter $B$ satisfies
\begin{equation}\label{eq:Bsolve}
B-\ln B=\ln\!\left(\frac{\chic^2}{1.167\,\chia^2}\right).
\end{equation}
For a common momentum through multiple materials, Lynch and Dahl's serial prescription gives
\begin{equation}\label{eq:multi}
\chic^2=\sum_i\frac{0.157\,Z_i(Z_i+1)X_i/A_i}{(p\beta)^2},
\end{equation}
\begin{equation}
\ln\chia^2=\frac{\sum_i w_i\ln\chi_{a,i}^2}{\sum_iw_i},\qquad
w_i=\frac{Z_i(Z_i+1)X_i}{A_i}.
\end{equation}
Equation~\eqref{eq:multi} is a common-$p$ relation; it should not be applied unchanged when the momentum varies appreciably along the path.

For the central Al+Cu reference path ($10~\mathrm{cm}$ Al, $27.0~\mathrm{g\,cm^{-2}}$, $\XZero=24.01~\mathrm{g\,cm^{-2}}$; $15~\mathrm{cm}$ Cu, $134.4~\mathrm{g\,cm^{-2}}$, $\XZero=12.86~\mathrm{g\,cm^{-2}}$) these give $(x/\XZero)_{\rm tot}=11.57$, $\sum_iw_i=2.02\times10^3$, $R=0.0621$, and $B=16.9$; for $15~\mathrm{cm}$ Pb, $R=0.0702$ and $B=16.5$.  These are the values quoted in Secs.~\ref{sec:RBinvariance} and~\ref{sec:collapse_results}.

\section{Asymptotic and truncation checks}\label{app:checks}
Figure~\ref{fig:modelchecks}(a) evaluates the absolute radial tail normalization through
\begin{equation}
\frac{h_{\rm M}(\Theta)\Theta^3}{2\chic^2}\longrightarrow1.
\end{equation}
For the $6~\mathrm{GeV}/c$ central Al+Cu path, the ratio is $1.349$, $1.073$, and $1.019$ at $50$, $100$, and $200~\mathrm{mrad}$, respectively.  The $200~\mathrm{mrad}$ point is therefore close to the Rutherford normalization, while $50~\mathrm{mrad}$ is visibly not yet in the asymptotic regime.  Panel (b) compares the accepted-moment mismatch at $n\le1$ and $n\le2$; the largest shift over the four momentum settings is about $0.062$ percentage point.

\begin{figure*}[t]
\centering
\includegraphics[width=0.48\textwidth]{figs/theory_tail_check}\hfill
\includegraphics[width=0.48\textwidth]{figs/theory_truncation_sensitivity}
\caption{Radial-model checks.  (a) Approach to the Rutherford tail normalization.  (b) Change in $\epsM$ when the Moli\`ere series is extended from $n\le1$ to $n\le2$.}
\label{fig:modelchecks}
\end{figure*}

The large-$\eta$ slope and intercept are checked separately in Fig.~\ref{fig:eta1}.  The slope diagnostic approaches the coefficient $2R$ as the fit window moves into the tail, with the deepest displayed window agreeing at approximately the percent level.  The fitted intercept $\eta_1$, however, is not universal: the asymptotic joint fits range from roughly $0.49$ for Pb to $1.35$ for Al, with Cu and Al+Cu near unity.  The $n\le1$ and $n\le2$ intercept estimates nearly overlap.  We therefore retain $\eta_1$ as a path-dependent core-to-tail matching parameter and do not identify it with a universal Moli\`ere scale.

\begin{figure*}[t]
\centering
\includegraphics[width=\textwidth]{figs/theory_eta1_paths}
\caption{Asymptotic-intercept diagnostic.  (a) Path-dependent fitted $\eta_1$ for $n\le1$ and $n\le2$.  (b) Fitted logarithmic slope relative to the exact tail coefficient $2R$ as the fit window moves to larger $\eta$.}
\label{fig:eta1}
\end{figure*}

\section{Segmented $p(X)$ construction and scope}\label{app:pofx}
The energy-loss-aware calculation is a controlled extension of the common-$p$ formulas, not a replacement for transport validation.  The mean energy is propagated through each material using the implemented Bethe collision stopping power with a Sternheimer density-effect correction.  Radiative loss is omitted.  Each ordered material segment is subdivided until $|\Delta(p\beta)/(p\beta)|<1\%$ per slice, with the representative $(p\beta)_j$ chosen to reproduce the numerical integral of $(p\beta)^{-2}$ over that slice.

The characteristic Moli\`ere strength is accumulated as
\begin{equation}\label{eq:chic_pofx}
\chic^2[p(X)]=\sum_j
\frac{0.157\,Z_j(Z_j+1)X_j/A_j}{(p_j\beta_j)^2}.
\end{equation}
The screening-log average is not uniquely specified by the common-$p$ serial formula once momentum varies along the path.  The implementation therefore retains two explicit continuations.  The production default weights each local $\ln\chi_{a,j}^2$ by the corresponding local increment $d\chi_{c,j}^2$; an alternative retains the common-$p$ serial material weight $Z_j(Z_j+1)X_j/A_j$ while evaluating $\chi_{a,j}$ at the local $p_j$ and $\beta_j$.  The two prescriptions are algebraically identical at constant momentum.  For the central Al+Cu path, their difference in the deployed upstream-tagged mismatch is $0.062$, $0.016$, $0.005$, and $0.002$ percentage point at $1$, $2$, $3.5$, and $6~\mathrm{GeV}/c$, respectively.  This spread is reported as an internal construction systematic; neither prescription is promoted to a transport theorem.

For the reference core, the construction used in the numerical study is
\begin{align}
\theta_{0,p(X)}^2
&=L_{\rm H}^2\sum_j
\left(\frac{13.6~\mathrm{MeV}}{p_j\beta_j c}\right)^2
\frac{\Delta x_j}{X_{0,j}},\label{eq:highland_pofx}\\
L_{\rm H}
&\equiv1+0.038\ln\!\left(\frac{(x/X_0)_{\rm tot}}{\beta_{\rm eff}^2}\right),
\end{align}
with $\beta_{\rm eff}^2$ weighted by the same local Gaussian scattering-power terms.  The empirical Highland logarithmic factor is applied once to the full path, not independently to each slice.  Equation~\eqref{eq:highland_pofx} reduces algebraically to Eq.~\eqref{eq:highland} when $p_j$ is constant, but it remains a construction choice.

The production geometry does not infer traversal order from unordered material totals.  The straight incoming ray is intersected analytically with the nested Al cube, Cu cube, and Pb cylinder, producing the ordered path $[\mathrm{Al}_{\rm up},\mathrm{Cu}_{\rm up},\mathrm{Pb},\mathrm{Cu}_{\rm down},\mathrm{Al}_{\rm down}]$ for every event.  In the normalization reference, the Pb interval is replaced by Cu.  This removes the earlier symmetric-order approximation for nonaxial rays.  The scattering cache then rounds incident momentum to $0.010~\mathrm{GeV}/c$, each ordered areal-density segment to $0.25~\mathrm{g\,cm^{-2}}$, and the cut to $0.002~\mathrm{rad}$.  A separate approximation remains at the detector-simulation level: after the integrated path distribution is sampled, the total deflection is represented by one equivalent kink at the midpoint of the outer-target traversal.

The fixed-cut results should therefore be read as a controlled $p(X)$ model study with exact ordered straight-ray geometry and verified constant-$p$ closure, not as a final detector-material systematic uncertainty.

\section{Independent transport benchmark protocol}\label{app:g4}
The reduced identity of Eq.~\eqref{eq:reduced_master} and the tail scaling of Eqs.~\eqref{eq:tail}--\eqref{eq:secondmoment} are internal statements of the analytic model.  Independent transport is therefore treated as a separate benchmark rather than as part of the production event generator.  Geant4 multiple-scattering implementations and the Urban/Wentzel model dependence have been documented and benchmarked against theory and data~\cite{agostinelli2003,allison2016,ivanchenko2010,makarova2017,mice2022}.

The supplied Geant4 executable fires a monoenergetic $\mu^-$ beam along the axis of a single Cu or Pb slab and records the primary exit-angle magnitude at a downstream scoring plane.  Each run takes an explicit random seed.  Three transport configurations are exposed: the unmodified \textsc{ftfp\_bert} reference list (``ftfp\_bert''), the unmodified \textsc{ftfp\_bert\_wvi} reference list (``ftfp\_bert\_wvi''), and an optional diagnostic configuration (``wvi\_ss'') that replaces the reference-list muon multiple-scattering process with a Wentzel-VI model plus discrete Coulomb scattering.  The reference-list names are retained because their installed muon models are Geant4-version dependent; in Geant4 11.4.2 both lists report WentzelVI for muons, with different settings.  The diagnostic configuration is not assumed a priori to be more physical than the reference lists.  The executable prints the installed muon process names at runtime so that every angle dump can be associated with its process configuration.

The Python comparison script accepts arbitrary single-material slab thicknesses and also the Al+Cu/Al-only reference paths for future layered transport runs.  At a cut $k=\thetacut/\thetazero$ it evaluates the accepted fraction, truncated RMS, and fourth moment.  The fractional RMS convention is fixed to the manuscript definition
\begin{equation}
D_{\rm RMS}\equiv\frac{\thetarms^{\rm model}}{\thetarms^{\rm G4}}-1.
\end{equation}
Because the tomography weight is quadratic in the inverse denominator, the script also reports the corresponding normalization bias
\begin{equation}
D_w=\left(\frac{\thetarms^{\rm G4}}{\thetarms^{\rm model}}\right)^2-1
=\frac{1}{(1+D_{\rm RMS})^2}-1.
\end{equation}
Thus an RMS-scale discrepancy should not be interpreted directly as an equal-sized error in the quadratic weight.

For a central-core diagnostic, the Geant4 radial median $m$ is converted to the Rayleigh core scale $s=m/\sqrt{\ln2}$ and compared with the model $\thetaspace$.  This is robust against rare extreme scatters but still tests the Rayleigh-core hypothesis only indirectly.  A central projected-angle fit over the Lynch--Dahl fraction would be a more direct check of the coefficient used in Eq.~\eqref{eq:highland}.

To localize a second-moment discrepancy, the comparison code additionally decomposes the unnormalized accepted second-moment numerator into finite reduced-angle bands.  For band $u\in[u_1,u_2]$, with $u=\Theta/\thetazero$, it records both the probability mass and
\begin{equation}
N_2(u_1,u_2)=\int_{u_1\thetazero}^{u_2\thetazero}
\Theta^2 h(\Theta)\,d\Theta .
\end{equation}
This directly distinguishes a core, shoulder, and far-tail contribution without inferring the location of a discrepancy from a single cumulative RMS.

No numerical Geant4 agreement table is asserted in this code-aligned source because the supplied source tree does not include the complete fresh angle dumps and run metadata needed to regenerate the previous transport numbers.  The benchmark tables and figures should be populated only from reruns that record the Geant4 version, transport mode, seed, number of generated primaries, number of recorded exits, and the corresponding comparator output.  The segmented $p(X)$ construction remains unvalidated until a layered or energy-loss-aware transport comparison is performed.

\section{Code availability}
The analysis and simulation code is available at \url{https://github.com/melora1/highland-mst}.  The submitted version should cite a tagged release and immutable archive corresponding exactly to the figures and tables used in the final manuscript.

\bibliographystyle{unsrt}
\bibliography{references_rev13}

\end{document}
